\subsection{Initial Research Objective}\label{route:problem}

The original goal was to improve the exact bilinear complexity of ordinary
$3\times3$ matrix multiplication: construct an algorithm with at most
22 products or establish the optimality of the known 23-product bound.
The goal concerns exact identities in the bilinear model. In this model
each counted product multiplies a linear form in one input by a linear
form in the other.

Strassen's seven-product $2\times2$ algorithm established the power of
bilinear reorganization~\cite{strassen1969gaussian}. Laderman supplied the
23-product $3\times3$ construction~\cite{laderman1976noncommutative};
Bl\"aser developed lower bounds for small
formats~\cite{blaser2003}. SAT-based algorithm discovery and AlphaTensor
show that machine search can discover exact multiplication
schemes~\cite{heule2019waysa,alphatensor2022}. A new candidate must satisfy all coefficient
identities, and a lower bound must cover all permitted decompositions.

Two broad directions developed alongside each other. One explored
finite-field restrictions, combinatorial supports, and certificates.
The other explored characteristic-zero deformations, rank-one-span
representations, and symmetry. Their relationship was not a contest
between fixed methods: each could supply examples, constraints, or
counterchecks for the other.

Figure~\ref{fig:research-route} summarizes these connections. Qiushi Engine
developed examples, formulations, and tests in both directions; their
mathematical consequences determined which questions remained productive.

\subsection{Tensor Conventions and Verification Controls}\label{route:controls}

Qiushi Engine first established a common mathematical convention.
Matrices were flattened in row-major order, the tensor was generated
from the multiplication rule, and known schemes were checked against
all 729 tensor entries. The same convention had to survive transposition,
cyclic permutation, quotient maps, and changes of basis.

This is not only implementation detail. A matrix representing a linear
functional and a matrix representing a vector transform contragrediently.
A quotient class is not generally a vector in an invariant coordinate
complement. Confusing either pair changes the encoded mathematical
problem while leaving arrays of plausible dimensions.

Known exact decompositions served as positive controls. Altered candidates
and deliberately inconsistent instances tested rejection. Such controls
can falsify an encoding quickly; passing them alone does not prove that
the encoding covers every legal decomposition. The final proof therefore
also requires an explicit map from a hypothetical decomposition to the
integer constraints and from those constraints to the checked formulas.

\subsection{Deformation, Pairing Defects, and the Limits of Numerical Evidence}
\label{route:deformation}

One natural construction route starts near a known rank-23 scheme.
The research examined tangent directions, second-order compatibility,
and continuation toward a rank-defective pairing. The motivation was to
identify a redundant product that could be removed while preserving the
target tensor.

The first important negative lesson concerned compactness. Several
apparently promising trajectories decreased a residual while the
factor amplitudes grew. Large terms cancelled rather than converging
to a finite exact rank-22 decomposition. Regularization exposed a
residual--amplitude tradeoff instead of establishing an exact solution.
The observed tradeoff identified a failure mechanism in these trajectories.
Other components of the exact solution space remained to be studied.

The representation was then revised. Pairing defects were formulated
as incidence conditions; full-span requirements and exceptional strata
were separated from numerical conditioning. Rank-one-span formulations
ask whether the nine-dimensional slice space lies inside a span generated
by the desired number of rank-one matrices. Completion charts attempted
to enforce this geometry more directly.

Several overextensions were corrected along the way. A tangent
calculation at a full-rank pairing is not automatically a tangent-cone
calculation on a rank-defect variety. A parameter count is not an
existence theorem. Local numerical failure, even around many known
schemes, does not control disconnected or exceptional components.
No characteristic-zero rank-22 construction or global rank-23
optimality proof resulted from this line.

\subsection{Quotient Cores and the Binary-Field Reframing}\label{route:cores}

Wang's certified lower-bound framework supplied a different entry point:
study what remains after removing directions from a factor
space~\cite{wangautomated}. The lower-bound-20 certificate was treated
as a fixed mathematical input, with its quotient semantics reconstructed
under the tensor convention used throughout the research.

The $E_{11}$ core provided a useful two-sided research question.
After removing that first-factor direction, the tensor has shape
$8\times9\times9$. Its rank-19 construction would embed back into the
full tensor, with three additional rank-one products restoring the
deleted slice. Thus a complete core witness would yield a full
rank-at-most-22 construction over $\FF$.

The converse was more delicate. Excluding one core family does not
automatically exclude every full 20-term decomposition. A lower-bound
argument must justify the reduction to that family, its symmetry
normalization, multiplicities after projection, and coverage of all
remaining branches. The research repeatedly returned to these
obligations rather than identifying a difficult subproblem with the
whole theorem.

The core had 255 nonzero first-factor directions and a complete
417,199-subspace occupation surface. These are finite, exact objects,
but their possible supports still form a large combinatorial space.
Reframing the goal over $\FF$ made exact tests available throughout
the search; identifying useful structure remained the central difficulty.

\subsection{Occupation Constraints and Factor Completion}\label{route:completion}

For a decomposition of length $n$, a certified quotient bound $L(W)$
implies
\[
 m(W)\le n-L(W).
\]
This inequality records how many first factors can occupy a subspace.
It deliberately forgets the second and third factors. A support
satisfying every such inequality can therefore remain unrealizable.

Qiushi Engine examined fixed-first-factor SAT formulations, shared
factor conditions from contractions, and lift-back problems from
smaller quotient tensors. A useful distinction emerged among three
levels:
\[
 \begin{gathered}
 \text{admissible projected support}\\
 \downarrow\\
 \text{compatible lift to the full factor space}\\
 \downarrow\\
 \text{complete bilinear decomposition}.
 \end{gathered}
\]
Each arrow needs additional conditions. Solving an earlier level is
not a witness at a later one.

Residual-capacity searches and learned violated rows reduced some
instances, but a constraint valid only after choosing a prefix could
not be reused globally without its conditioning hypotheses. Quotient
lift-bit formulations also required the same coordinate map in the
row exporter and the branch solver. A mismatch once produced an
apparent exclusion that was withdrawn when the semantics were
reconstructed.

Repeated projected directions formed a separately treated branch;
the distinct-direction branch remained difficult. The eventual
full-tensor proof did not require the exact rank of the entire
$E_{11}$ core to be determined.

\subsection{Symmetry-Restricted Constructions and Their Scope}
\label{route:symmetry}

The characteristic-zero line explored inner and outer tensor symmetries,
cyclic decompositions, transpose-invariant families, and mixed
symmetry types. Exterior and symmetric components, residual cubics,
commutator or Pfaffian tests, and rank-one-span conditions offered
smaller algebraic descriptions than unrestricted Brent equations.

This line generated local obstructions and exact test instruments.
Its main caution was equally important: a result for a chosen frame,
seed orbit, or symmetry skeleton cannot be promoted to an unrestricted
rank bound. Finite-field reductions must retain their field scope;
changes of frame may preserve one exterior component while changing
the residual symmetric component that a lower-bound test consumes.

Several apparent shortcuts did not survive scrutiny. A direct-sum
rank inference was too broad. Some tests used a dual action in the
wrong coordinates. A proposed restriction on a transpose family
failed when compared with a correctly transported known construction.
The corrected calculations distinguished invariants of the complete
tensor from properties of one chosen representation.

These explorations supplied local algebraic descriptions and tests
for symmetry-based construction. The binary-field proof ultimately
followed a different route, using constraints on every hypothetical
20-term decomposition rather than a selected symmetry family.

\subsection{Relaxations and Information Loss}\label{route:relaxation}

The difficult distinct-support core branch motivated coding-theoretic,
linear, pair-correlation, and semidefinite relaxations. Orbit-averaged
pair variables and moment matrices compressed large labelled systems.
Exact square inequalities could sometimes turn a numerical suggestion
into a reusable, auditable condition.

These experiments clarified a hierarchy of evidence. A fractional
occupation vector is not an integer support. An integer pair-count
profile need not arise from a labelled point set. A positive
semidefinite numerical matrix does not itself provide the required
factor decomposition. Conversely, an ill-conditioned conic solve
does not establish exact infeasibility.

Some intermediate conclusions had to be narrowed after missing
triangle constraints, forced kernels, normalization factors, or
different row families were identified. A cut that separates one
saved fractional point does not eliminate all feasible supports.
The sustained lesson was that each compression must state what
information it retains and what remains to be reconstructed.

No complete core exclusion emerged from these relaxations. The
research therefore kept their useful constraints without treating
ever larger relaxations as the only possible next step.

\subsection{Full-Tensor Rank Constraints}\label{route:full}

The split flattening reorganizes the full multiplication tensor into
a $27\times27$ permutation matrix $P$. A simple tensor
$A_t\otimes B_t\otimes C_t$ becomes a matrix $M_t$ of rank
$\rk(A_t)$ when the other two factors are nonzero. Consequently,
\[
 27=\rk(P)\le\sum_{t=1}^{n}\rk(A_t).
\]
At $n=20$, the excess above twenty rank-one first factors must total
at least seven. At least four first factors must therefore have
matrix rank at least two.

This bound supplied information that first-factor occupation alone
had not expressed sharply. It also suggested a more economical
question: which low-dimensional quotient raises constrain pairs
of high-rank factors strongly enough to determine their geometry?

The eventual proof needs eight two-dimensional quotient orbits.
Each consists of planes whose three nonzero matrices all have rank
at least two. Raising their quotient lower bounds from 18 to 19
means that two high-rank first factors cannot span such a plane
in a 20-term decomposition. Their sum must therefore have rank one.

\subsection{Affine Geometry and the Forced Rank Profile}\label{route:geometry}

A family of matrices with pairwise rank-one differences lies in a
common affine row or column coset, by classical matrix
geometry~\cite{hua1945geometries,wan1996geometry}. This identification
turned the pairwise rank condition into a small affine configuration
whose occupation could be analyzed explicitly.

In the relevant normalized coset, eight first-row labels form
$\FF^3$. Four corresponding matrices have rank two and four have
rank three. Certified affine-plane caps permit at most three chosen
points on any four-point affine plane. Every five-point subset of
$\FF^3$ contains such a plane, so at most four high-rank factors
can occur.

Flattening demands at least four. Among four, all four rank-three
points would themselves form a forbidden affine plane. To supply
the required rank excess, the only remaining possibility is
one rank-two and three rank-three factors. The complete profile is
\[
 (n_1,n_2,n_3)=(16,1,3),\qquad
 n_1+2n_2+3n_3=27.
\]
This equality was the decisive interface between combinatorial
geometry and algebraic rigidity.

\subsection{Saturation and Cross-Factor Rigidity}\label{route:saturation}

For an invertible matrix $P=\sum_t M_t$, equality
$\sum_t\rk(M_t)=\rk(P)$ makes the column spaces a direct sum.
Resolving a vector in that direct sum gives
\[
 M_tP^{-1}M_s=\delta_{ts}M_t.
\]
For the multiplication tensor, direct index contraction yields
\[
 M(A,B,C)P^{-1}M(A',B',C')
   =M\bigl(A(B'C^T)A',B,C'\bigr).
\]
The diagonal identity for an invertible $A_t$ forces
$B_tC_t^T=A_t^{-1}$, so $B_t$ and $C_t$ are invertible.
If two first factors $A_t,A_s$ are invertible, the off-diagonal
identity requires $A_tB_sC_t^TA_s=0$, impossible for a product
of invertible matrices.

Thus a saturated decomposition permits at most one invertible
first factor. The affine profile requires three. This is the
contradiction proving the lower bound 21.

Qiushi Engine had changed the final obligation: a potentially large
terminal enumeration became an exact cross-factor identity. An unnecessary
additional quotient orbit was also removed from the assumptions.
The final argument retains the eight essential finite raises
and eliminates the larger terminal search from its dependency chain.

\subsection{Certification of the Finite Premises}\label{route:certificates}

Each essential quotient has 127 nonzero directions and 29,210
nonzero proper subspaces. At target length 18, its occupation
variables are nonnegative integers, constrained to sum to 18.
Their capacities come from the verified upstream quotient table.

Two encoding approaches were developed. Bounded Boolean copies
represent each integer multiplicity as a sum of binary variables.
Monotone unary counts represent whether a multiplicity reaches
each possible level. The latter exposes the count semantics
directly. Both must allow multiplicity two when the quotient
permits it; replacing multiplicities by distinct-support bits
would change the mathematical question.

An earlier cardinality encoding accidentally reused auxiliary
variables. The resulting contradiction was an error in the
formula, not an exclusion of a tensor decomposition. Its apparent
lower-bound conclusion was withdrawn. The original SAT verification
boundary therefore has two parts:
\[
 \text{correct mathematical encoding}
 \qquad+\qquad
 \text{correct proof of formula unsatisfiability}.
\]
DRAT establishes the second part. Explicit semantic reconstruction,
positive controls, and separate encoders address the first.

All eight final clausal proofs are retained. Their formulas can
be regenerated and compared byte for byte. The source certificate
can be reverified, and its orbit representatives independently
expanded to the complete 8,283,458-subspace table. The structural
argument then derives the contradiction from these finite premises.

The accompanying Lean development proves the finite premises by integer
branch certificates on the actual quotient objects and connects them to
the structural contradiction. It also checks the orbit classification,
occupation systems, calibration controls, and Boolean encoding semantics.
These formal counterparts preserve the mathematical lessons of the
computational work while making the complete theorem a kernel-checked
dependency chain (Section~\ref{sec:lean}).

\subsection{Corrections and Methodological Consequences}\label{route:corrections}

\begin{longtable}{@{}L{3.3cm}L{5.1cm}L{5.8cm}@{}}
\toprule
Issue & Why the apparent inference failed & Retained scientific consequence\\
\midrule
\endfirsthead
\toprule Issue & Why the apparent inference failed & Retained scientific consequence\\\midrule
\endhead
\bottomrule\endfoot
Near-zero numerical residual & Coefficients could diverge while cancelling
& Require a finite exact identity; distinguish tested degeneration from impossibility\\
Auxiliary-variable collision & The Boolean formula contradicted itself for encoding reasons
& Withdraw the claim; verify variable allocation and semantic forward maps\\
Lookup decoding & A packed subspace key was read with the wrong bit convention
& Reconstruct canonical row spaces and compare the entire expanded table\\
Quotient coordinates & A projection and a lift used different coordinate maps
& Keep the quotient map, complement, and induced action explicit\\
Local symmetry exclusion & A seed, orbit, or frame did not cover all decompositions
& State exact hypotheses and retain uncovered branches\\
Relaxation feasibility & Fractional or averaged moments need not realize a discrete support
& Separate continuous, integral, labelled, and factor-complete objects\\
Product indexing & An elementary formula did not match the actual permutation contraction
& Construct the permutation matrix independently and check the elementary products\\
Unnecessary finite premise & One extra pair orbit already had rank-one differences
& Remove it from the shortest proof instead of proving a stronger statement unnecessarily\\
\end{longtable}

The useful outcome of correction was a change in subsequent practice.
The cardinality error led to explicit assignment maps and separate
count encodings. Coordinate mismatches led to reconstruction from the
tensor and its maps. Discovering that one finite premise was unnecessary
led to a shorter argument. In each case the retained knowledge was more
specific than an instruction to check carefully: it identified the
mathematical relation that later computations had to preserve.

\subsection{Results, Reusable Methods, and Open Problems}\label{route:outcome}

The established result is
\[
 R_{\FF}(\Tmmm)\ge21.
\]
The final proof combines certified quotient bounds, affine rank-one
geometry, rank-additive split flattening, and a tensor-specific product
identity. It also shows that a zero-excess rank-22 decomposition in
one factor slot can have at most one invertible factor in that slot.
When all 22 factors in that slot are nonzero, the surviving rank profiles
are $(17,5,0)$ and $(18,3,1)$. With zero factors allowed, the alternative
reduction to a zero-excess length-21 decomposition is retained as in
Section~\ref{sec:outlook}.

The lower bound and these structural consequences have completed Lean
proofs within the hypotheses recorded in \qpath{formalization/COVERAGE.md}.
The main theorem has no remaining finite-bound assumptions.

The exact binary-field rank remains in $[21,23]$. Constructing a
rank-22 algorithm, controlling positive-excess decompositions, and
extending the mechanism to other fields are natural next questions.
The quotient and symmetry materials preserve routes to those questions
alongside the completed lower-bound proof.

The Meta-Trace records a second outcome: the development of a research
method through accumulated understanding. Qiushi Engine constructed
representations, used counterexamples and calculation to identify their
limits, and carried the useful distinctions into later reasoning. The
separation between support and completion became a question about how
to recover factor coupling; the forced equality supplied the answer.

Documenting this development together with its mathematical artifacts
allows more than replay of the accepted proof. Readers can inspect
the hypotheses behind an experiment, understand why its outcome changed
the route, adapt an intermediate method, or revisit a remaining question
with its earlier obstacles made explicit. The proof establishes the
result; the Meta-Trace preserves the reasoning history from which
further research can proceed.
